\section{Метрики устойчивости}

Для оценки физической устойчивости укладки реализованы метрики, предложенные в работе \cite{TR-MFTI-khelvas2019-boxes}.

% TODO: Добавить иллюстрацию хорошей и плохой перевязки (figures/stability_comparison.png)
% Показать: два примера укладки - слева "кирпичная кладка" (stability=0.9),
% справа "столбики" (stability=0.3), с выделенными совпадающими рёбрами

\subsection{Коэффициент stability}

Основной метрикой устойчивости является коэффициент \textbf{stability}, который вычисляется по формуле:

$$stability = 1 - InterlockingRatio = 1 - \frac{L_{Sum}}{L_{SumPer}}$$

\subsubsection{Состав метрики}

\begin{itemize}
    \item $L_{SumPer}$ --- суммарный периметр оснований \textbf{всех} коробок в укладке
    \item $L_{Sum}$ --- суммарная длина \textbf{совпадающих} кусков периметров коробок, стоящих друг на друге
    \item $InterlockingRatio = \frac{L_{Sum}}{L_{SumPer}}$ --- доля совпадающих рёбер
\end{itemize}

\subsubsection{Диапазон значений}

\begin{itemize}
    \item $stability \in [0, 1]$
    \item $stability = 1$ --- идеальная перевязка: ни одно ребро верхней коробки не совпадает с ребром нижней (как в кирпичной кладке)
    \item $stability = 0$ --- все рёбра совпадают: коробки стоят ``столбиками'' и могут рассыпаться
\end{itemize}

\subsubsection{Интерпретация}

\begin{table}[ht]
\centering
\caption{Интерпретация значений stability}
\begin{tabular}{|c|l|}
\hline
\textbf{Значение} & \textbf{Качество перевязки} \\
\hline
0.9 -- 1.0 & Отличная перевязка \\
0.7 -- 0.9 & Хорошая перевязка \\
0.5 -- 0.7 & Удовлетворительная \\
0.3 -- 0.5 & Слабая перевязка \\
0.0 -- 0.3 & Плохая перевязка (риск рассыпания) \\
\hline
\end{tabular}
\end{table}

\subsubsection{Ограничения метрики}

Метрика stability имеет ряд ограничений:

\begin{enumerate}
    \item \textbf{Не учитывает опорную площадь}: коробка может ``висеть в воздухе'' (опираться менее чем на 50\% площади), но метрика этого не отражает.
    
    \item \textbf{Не учитывает вес}: тяжёлая коробка на лёгкой может перевесить и упасть, даже если перевязка хорошая.
    
    \item \textbf{Не учитывает высоту}: коробки на большой высоте более чувствительны к смещению.
    
    \item \textbf{Упрощённая модель совпадения}: считаются только точные совпадения координат рёбер. Если рёбра отличаются на 1 мм, они не считаются совпадающими.
    
    \item \textbf{Не учитывает направление щелей}: вертикальная щель через всю паллету опаснее нескольких коротких.
\end{enumerate}

\subsubsection{Практическое применение}

Несмотря на ограничения, метрика полезна для:
\begin{itemize}
    \item Сравнения разных укладок одного набора коробок
    \item Быстрой оценки качества алгоритма
    \item Оптимизации: добавление штрафа за низкий stability в целевую функцию
\end{itemize}

\subsection{Метрика interlocking\_ratio}

\textbf{interlocking\_ratio} --- это доля совпадающих рёбер между соседними по высоте коробками:

$$interlocking\_ratio = \frac{L_{Sum}}{L_{SumPer}}$$

\subsubsection{Что измеряет}

Метрика показывает, какая часть периметров коробок ``совпадает'' с периметрами коробок, на которых они стоят.

\begin{itemize}
    \item $L_{Sum}$ --- считаем все пары коробок (A, B), где A стоит на B (то есть нижняя грань A касается верхней грани B). Для каждой такой пары находим длину отрезков, где рёбра коробок точно совпадают.
    \item $L_{SumPer}$ --- суммарный периметр оснований всех коробок.
\end{itemize}

\subsubsection{Диапазон значений}

\begin{itemize}
    \item $interlocking\_ratio \in [0, 1]$
    \item $interlocking\_ratio = 0$ --- рёбра не совпадают, идеальная ``кирпичная кладка''
    \item $interlocking\_ratio = 1$ --- все рёбра совпадают, коробки стоят ``столбиками''
\end{itemize}

\subsubsection{Связь со stability}

$$stability = 1 - interlocking\_ratio$$

То есть: \textbf{чем меньше interlocking\_ratio, тем лучше} (меньше совпадающих щелей).

\subsection{Метрика center\_of\_mass\_z}

\textbf{center\_of\_mass\_z} --- высота центра масс паллеты:

$$center\_of\_mass\_z = \frac{\sum_{i=1}^{K} m_i \cdot z_i^{center}}{\sum_{i=1}^{K} m_i}$$

где:
\begin{itemize}
    \item $m_i$ --- масса $i$-й коробки (поле Weight из входных данных)
    \item $z_i^{center} = \frac{z_i + Z_i}{2}$ --- высота центра $i$-й коробки
\end{itemize}

\subsubsection{Что измеряет}

Показывает, насколько высоко расположен центр тяжести всей паллеты. Единица измерения --- миллиметры (как и все координаты).

\subsubsection{Диапазон значений}

\begin{itemize}
    \item $center\_of\_mass\_z \in [h_{min}/2, H_{total} - h_{min}/2]$
    \item Типичные значения: от 200 до 600 мм (зависит от высоты паллеты)
\end{itemize}

\subsubsection{Интерпретация}

\begin{itemize}
    \item \textbf{Низкий центр тяжести} (близко к 0) --- паллета устойчива, тяжёлые коробки внизу
    \item \textbf{Высокий центр тяжести} (близко к $H_{total}/2$) --- паллета менее устойчива, может опрокинуться при транспортировке
\end{itemize}

\subsubsection{Ограничения}

\begin{enumerate}
    \item Если вес коробок не указан (Weight = 0), используется вес = 1 для всех коробок
    \item Не учитывается смещение центра тяжести по X и Y (только по высоте)
    \item Не учитывается форма паллеты и погрузчика
\end{enumerate}

\subsection{Реализация}

Метрики устойчивости реализованы в модуле \texttt{stability.cpp}:

\begin{itemize}
    \item \texttt{calc\_stability()} --- рассчитывает все метрики для заданной укладки
    \item \texttt{calc\_overlapping\_edges()} --- определяет длину совпадающих рёбер между парой коробок
    \item \texttt{segment\_overlap\_length()} --- вычисляет длину пересечения двух отрезков
\end{itemize}

Результаты выводятся в файл \texttt{metrics.csv} с колонками:
\begin{verbatim}
stability, interlocking_ratio, center_of_mass_z
\end{verbatim}
